% Part: normal-modal-logic
% Chapter: filtrations
% Section: finite

\documentclass[../../../include/open-logic-section]{subfiles}

\begin{document}

\olfileid[pt]{nml}{fil}{fin}

\olsection{As filtrações são finitas}

Definimos filtrações para qualquer conjunto~$\Gamma$ fechado sob
sub-!!{formula}s. Nada na própria definição garante que as filtrações
sejam finitas. De fato, quando $\Gamma$ é infinito (por exemplo, quando é
o conjunto de todas as !!{formula}s), a filtração pode muito bem ser
infinita. Contudo, se $\Gamma$ é finito (por exemplo, quando é o conjunto
das sub-!!{formula}s de uma dada !!{formula}~$!A$), então toda filtração
por meio de~$\Gamma$ também é finita.

\begin{prop}\ollabel{prop:filt-are-finite}
  Se $\Gamma$ é finito, então toda filtração $\mModel{M^*}$ de um modelo
  $\mModel{M}$ por meio de $\Gamma$ também é finita.
\end{prop}

\begin{proof}
  O tamanho de $W^*$ é o número de classes distintas~$[w]$ sob a relação
  de equivalência~$\equiv$. Quaisquer dois mundos $u$, $v$ em uma mesma
  classe---isto é, quaisquer $u$ e $v$ tais que~$u \equiv v$---concordam
  sobre todas as !!{formula}s~$!A$ em~$\Gamma$: para $!A \in \Gamma$,
  ou $!A$ é verdadeira tanto em $u$ quanto em $v$, ou não é verdadeira
  em nenhum deles. Assim, cada classe~$[w]$ corresponde a um subconjunto
  de~$\Gamma$, a saber, o conjunto de todos os $!A \in \Gamma$ tais que
  $!A$ é verdadeira nos mundos de~$[w]$. Duas classes diferentes $[u]$
  e $[v]$ não correspondem ao mesmo subconjunto de~$\Gamma$. Com efeito,
  se o conjunto das !!{formula}s verdadeiras em $u$ é o mesmo que o das
  !!{formula}s verdadeiras em $v$, então $u$ e $v$ concordam sobre todas
  as fórmulas em~$\Gamma$, isto é, $u \equiv v$. Mas, nesse caso,
  $[u] = [v]$. Logo, há uma função !!a{injective} de $W^*$ em
  $\Pow{\Gamma}$ e, portanto, $\card{W^*} \le \card{\Pow{\Gamma}}$.
  Assim, se $\Gamma$ contém $n$ sentenças, a cardinalidade de $W^*$ não
  é maior que~$2^n$.
\end{proof}

\end{document}


